

\documentclass[tikz,svgnames,border={10 10},convert={density=300, outext=.png}]{standalone}

% 若需要中文，请用 xelatex 编译并取消下面两行注释：
\usepackage{xeCJK}
% \setCJKmainfont{Noto Sans CJK SC}

\usetikzlibrary{positioning,arrows,fit}

\begin{document}
\begin{tikzpicture}[
    box/.style={rectangle,draw,fill=DarkGray!20,rounded corners,
                node distance=1cm,text width=18em,text centered,
                minimum height=2.4em,thick},
    arrow/.style={draw,-latex',thick},
  ]

  % === 主体流程节点（与前述模块一致风格，不加外框） ===
  \node[box,fill=PeachPuff!60,text width=24em] (input)
    {接收数据源与请求：CSV/Parquet/数据库、实时流、合成参数；\ 外部调用 load\_dataset / get\_snapshot / request\_data / subscribe};

  \node[box,below=0.6 of input] (load)
    {历史/实时数据加载（load\_dataset；mmap/Parquet/Kafka/WebSocket）};

  \node[box,below=0.5 of load] (cache)
    {缓存与索引维护（LRU / 时间索引 / 权限校验 / 时间对齐）};

  \node[box,below=0.5 of cache] (route)
    {查询与订阅路由（get\_snapshot / request\_data / subscribe）};

  \node[box,below=0.5 of route] (crit)
    {数据就绪并满足请求？};

  % 底部输出
  \node[box,below=1.8 of crit,fill=LightSteelBlue!60,text width=26em] (output)
    {输出：时间对齐快照、tick / OHLC、Future / 回调；\ 持久化日志与元信息记录};

  % === 箭头（直连 + Yes/No 回路） ===
  \path[arrow] (input) -- (load);
  \path[arrow] (load) -- (cache);
  \path[arrow] (cache) -- (route);
  \path[arrow] (route) -- (crit);

  % Yes 分支：直达输出
  \path[arrow] (crit) -- (output)
    node[midway,left=0.1,draw,outer sep=0pt,fill=white] {Yes};

  % No 分支：先竖直向下，再水平右移（No 在水平段下方），最后折返到“查询与订阅路由”节点
  \path[draw,thick] (crit.south) ++(0,-1em) coordinate(novert);
  \draw[thick] (novert) -- ++(4,0) coordinate(nojump)
        node[midway,below=2pt,draw,outer sep=0pt,fill=white] {No};
  \draw[arrow] (nojump) |- (route.east);

\end{tikzpicture}

\end{document}